\section{Methodology}


We describe the proposed GAD-MambaUNet by first introducing its core building block and then integrating it into an asymmetric U-shaped segmentation architecture. 
As shown in Fig.~\ref{fig:overall-network}, the left part presents the overall student network and the training-time DINOv3 teacher, while the right part details the proposed block-level design and the internal structure of Direction-Group Graph Selective Scan (DG-GSS). 
We first present DG-GSS, which enables structured interaction among scan directions and channel groups. Then, we describe how DG-GSS is embedded into the overall GAD-MambaUNet architecture. Finally, we introduce the training-time DINOv3 supervision with Gradient-Adaptive Distillation (GAD), where the teacher model is used only during training and removed during inference.



% \section{Methodology}

% \subsection{Overview}

% Fig.~\ref{fig:overall-network} presents GAD-MambaUNet. Given an image $\mathbf{x}\in\mathbb{R}^{3\times H\times W}$, the student produces a binary logit map $\mathbf{z}\in\mathbb{R}^{1\times H\times W}$. The base model uses five resolution levels with channel widths $[16,32,64,96,160]$. The first three encoder blocks and the last three decoder blocks are multi-kernel inverted residual (MKIR) blocks. The two deepest encoder blocks and the first two decoder blocks are grouped Mamba blocks. Resolution and channel changes are performed outside these blocks by depth-wise downsampling followed by point-wise projection, or by point-wise projection followed by bilinear upsampling. This separation keeps the state-space width fixed within each block.

% At every decoder level, channel and spatial attention refine the decoder feature, while a grouped attention gate filters the corresponding skip connection before additive fusion. The first decoder block operates at $H/32\times W/32$ and supplies the student feature used for DINOv3 supervision. This location is semantically rich, inexpensive, and directly aligned with the block monitored by GAD.

% \subsection{Asymmetric Local--Global Student}

% At high resolution, the dominant requirement is local boundary reconstruction. We therefore retain the MKIR design of MK-UNet~\cite{rahman2025mkunet}: a $1\times1$ expansion, parallel depth-wise convolutions with kernels $\mathcal{K}=\{1,3,5\}$, summation, channel shuffle, and a $1\times1$ projection. For feature $\mathbf{x}$, the block is written compactly as
% \begin{equation}
% \operatorname{MKIR}(\mathbf{x})=
% P_2\!\left(\operatorname{Shuffle}\!\left[
% \sum_{k\in\mathcal{K}}D_k\!\left(P_1(\mathbf{x})\right)\right]\right),
% \end{equation}
% where $P_1$ and $P_2$ denote point-wise expansion and projection and $D_k$ is a depth-wise convolution with kernel size $k$. The block supplies multi-receptive-field local evidence at low cost.

% At low resolution, the receptive field can be enlarged more efficiently. A grouped Mamba block first applies depth-wise local mixing and a grouped shuffle feed-forward network (FFN), then DG-GSS, and finally a second local mixing--FFN pair. Each component is residual:
% \begin{equation}
% \begin{aligned}
% \mathbf{u}_1 &= \mathbf{x}+\operatorname{FFN}_0(\operatorname{DW}_0(\mathbf{x})),\\
% \mathbf{u}_2 &= \mathbf{u}_1+\operatorname{DG\text{-}GSS}(\mathbf{u}_1),\\
% \mathbf{y} &= \mathbf{u}_2+\operatorname{FFN}_1(\operatorname{DW}_1(\mathbf{u}_2)).
% \end{aligned}
% \end{equation}
% The local operators therefore refine state-space features instead of competing with them.

% For skip fusion, let $\mathbf{g}$ denote the upsampled decoder feature and $\mathbf{s}$ the encoder feature at the same resolution. A grouped attention gate computes
% \begin{equation}
% \boldsymbol{\alpha}=\sigma\!\left(\psi\!\left(\phi(W_g\mathbf{g}+W_s\mathbf{s})\right)\right),\qquad
% \hat{\mathbf{s}}=\boldsymbol{\alpha}\odot\mathbf{s},
% \end{equation}
% where $W_g$ and $W_s$ are grouped convolutions, $\phi$ is ReLU, and $\psi$ maps the joint response to a spatial gate. The decoder feature and $\hat{\mathbf{s}}$ are then added.

% \subsection{Direction-Group Graph Selective Scan}

% Let $\mathbf{X}\in\mathbb{R}^{B\times D\times h\times w}$ be the input to DG-GSS. We use four Cross2D traversal directions and partition the inner state-space channels into $M$ groups. With $d=D/M$ and $L=hw$, cross scanning yields
% \begin{equation}
% \mathbf{X}_{\mathrm{scan}}\in\mathbb{R}^{B\times K\times M\times d\times L},\qquad K=4.
% \end{equation}
% Every direction--group pair independently produces its selective-scan parameters $(\Delta,\mathbf{B},\mathbf{C})$. Selective scan is applied to each pair and the four outputs are spatially realigned to the original coordinate system:
% \begin{equation}
% \mathbf{Y}\in\mathbb{R}^{B\times K\times M\times d\times h\times w}.
% \end{equation}

% The tensor $\mathbf{Y}_{k,m}$ is a response from direction $k$ and channel group $m$. We regard it as a graph node and obtain a node descriptor by global average pooling:
% \begin{equation}
% \mathbf{r}_{k,m}=\frac{1}{hw}\sum_{u,v}\mathbf{Y}_{k,m,:,u,v}.
% \end{equation}
% After layer normalization and linear projection, additive graph-attention logits~\cite{velickovic2018gat} are
% \begin{equation}
% e_{ij}=\operatorname{LeakyReLU}\left(\mathbf{a}_{s}^{\top}\mathbf{h}_i+
% \mathbf{a}_{t}^{\top}\mathbf{h}_j\right).
% \end{equation}
% Our factorized adjacency connects two nodes when they have the same scan direction or the same channel-group identity, including self-connections. Thus, a node receives both cross-group and cross-direction evidence without constructing an unconstrained fully connected graph. With masked softmax weights $a_{ij}$ and shared $1\times1$ value projection $V(\cdot)$, graph propagation is
% \begin{equation}
% \widetilde{\mathbf{Y}}_i=\mathbf{Y}_i+\gamma\sum_{j\in\mathcal{N}(i)}a_{ij}V(\mathbf{Y}_j),
% \end{equation}
% where $\gamma$ is learnable and initialized to zero. DG-GSS consequently starts from the original grouped scan behavior and learns graph exchange progressively. Finally, groups are concatenated and the four aligned directions are summed before output normalization.

% \subsection{DINOv3 Supervision and Gradient-Adaptive Distillation}

% A frozen DINOv3 ViT-B/16 teacher provides dense features only during training. The feature $\mathbf{F}_s$ after the first decoder block is projected by a $1\times1$ convolution and batch normalization to the teacher dimension. The teacher feature $\mathbf{F}_t$ is bilinearly resized when necessary. After flattening each map into spatial tokens and applying $\ell_2$ normalization, the feature loss is
% \begin{equation}
% \mathcal{L}_{\mathrm{dist}}=
% \frac{1}{N}\sum_{n=1}^{N}\left[1-
% \frac{\mathbf{f}_{s,n}^{\top}\mathbf{f}_{t,n}}
% {\|\mathbf{f}_{s,n}\|_2\|\mathbf{f}_{t,n}\|_2}\right].
% \end{equation}
% The segmentation objective follows the structure-aware weighted BCE and weighted IoU loss~\cite{fan2020pranet}:
% \begin{equation}
% \mathcal{L}=\mathcal{L}_{\mathrm{seg}}+\lambda_e\mathcal{L}_{\mathrm{dist}},
% \qquad
% \mathcal{L}_{\mathrm{seg}}=\mathcal{L}_{\mathrm{wbce}}+\mathcal{L}_{\mathrm{wiou}}.
% \end{equation}

% Using a fixed $\lambda_e$ assumes that the appropriate teacher pressure remains unchanged throughout training. Instead, after backpropagation and before gradient clipping, GAD measures the gradient share of the aligned decoder block:
% \begin{equation}
% q_e=100\times
% \frac{\sum_{\theta_i\in\Theta_a}\|\nabla_{\theta_i}\mathcal{L}\|_1}
% {\sum_{\theta_j\in\Theta}\|\nabla_{\theta_j}\mathcal{L}\|_1}.
% \end{equation}
% Here, $\Theta_a$ is the parameter set of decoder1 and $\Theta$ contains all student parameters. If $q_e$ falls outside the target interval $[\rho-\delta,\rho+\delta]$, GAD selects the opposite interval boundary as $q^*$ and computes an odds-ratio adjustment:
% \begin{equation}
% r_e=\frac{p^{*}(1-p_e)}{p_e(1-p^{*})},
% \quad p_e=q_e/100,\quad p^{*}=q^{*}/100.
% \end{equation}
% The next base distillation coefficient is
% \begin{equation}
% \lambda_{e+1}=\operatorname{clip}\left(
% \lambda_e\operatorname{clip}(r_e,r_{\min},r_{\max}),
% \lambda_{\min},\lambda_{\max}\right).
% \end{equation}
% The multiplicative adjustment is bounded per epoch and the coefficient is clipped to $[\lambda_{\min},\lambda_{\max}]$ for stability. Training begins with segmentation-only burn-in, follows with a linear distillation warm-up, and activates GAD after warm-up. At inference, the teacher, feature projection head, and GAD update are discarded.
